\chapter*{Introducción}

\textbf{Esta introducción es para la revisión del manuscrito, la cambiaré al final.}

\section*{Idea de tesis}

\begin{itemize}
    \item la idea de este trabajo de tesis es construir un camino natural entre los modelos de difusión y el problema del puente de Schrödinger de forma simple, sin entrar en tecnicismos que usualmente se ven en la mayoría de papers que he revisado del SBP. Pensar este trabajo como un documento inclinado más hacia la comunidad del aprendizaje automático que hacia la matemática en sí. Por esto mismo, algunas cosas no son del todo correctas (p.g. la hipótesis de medibilidad de las funciones por lo general es omitida y tampoco se entra en detalles como verificar la condición de Slater en los problemas de optimización). De todos modos, he intentado mantener lo más posible la rigurosidad en la presentación.
    \item Los modelos de difusión funcionan, constituyen el estado del arte y, en la práctica, generan resultados casi perfectos. Además, se usan para muchas cosas (generación de imágenes, videos, grafos, análisis molecular, sonido, medicina, etc.).
    \item Tienen limitaciones: hay que correr el proceso hasta un tiempo muy largo para intentar aproximar bien la distribución invariante (gaussiana estándar). esta inexactitud inhabilita los resultados teóricos que pueda haber, depende mucho del proceso de difusión elegido y hace que el proceso de generación sea muy lento. Por otra parte, si se intenta ver como un modelo que transforma distribuciones, tiene la limitación de que solo se puede llegar a una distribución gaussiana. Esto último motiva a ver si es posible estudiar los DM desde la perspectiva del transporte óptimo.
    \item Para solucionar los problemas mencionados, se propone estudiar el problema del puente de Schrödinger que transforma una distribución en otra (esta última arbitraria) en un horizonte de tiempo finito.
    \item Para estudiar este problema, se hará desde la perspectiva del transporte óptimo ya que es un área bien estudiada y tiene propiedades importantes. Más aún, se prueba en el capítulo 3 que la forma regularizada del transporte óptimo es equivalente al problema de Schrödinger, por lo que todo lo visto en el capítulo 2 se hereda (de forma aproximada) al problema de Schrödinger.
    \item En consecuencia, en el capítulo 2 se introduce el problema de Monge, se ven sus problemas y luego se introduce el problema de Kantorovich. Se estudian varias propiedades potentes de esta formulación y luego se muestra que, si bien es una formulación interesante, tiene limitaciones técnicas (en particular, es muy costoso e inestable). Se termina el capítulo concluyendo que los DM son (aproximadamente) un caso particular de transporte óptimo, pero para esto último es necesario haber estudiado en detalle la formulación continua de los DM, lo que justifica todo el trabajo realizado en el capítulo 1.
    \item En el capítulo 3 (incompleto, falta traspasar varias cosas), se estudia la forma regularizada del transporte óptimo y se estudian algunas de sus formulaciones dinámicas, las cuales son más o menos análogas a las del transporte óptimo no regularizado. Estas formulaciones, si bien parecen complejas, han sido esenciales para resolver el problema de Schrödinger desde la perspectiva del machine learning (en particular, las cosas difíciles las aproximan, como es usual, por redes neuronales).
    \item Es importante decir que, si bien la mayor parte del trabajo estudia la formulación estática de las cosas (i.e., buscar una función que transporte una medida en otra), la formulación dinámica (i.e., las cosas como procesos estocásticos) es directa ya que solo hace falta introducir un puente browniano entremedio de la formulación estática, lo cual es muy fácil de hacer.
    \item Por otra parte, la idea es que todo lo que hay en el documento sea útil (i.e., no rellenar con texto innecesario), por lo que si sobra algo se puede eliminar. Además, todas las ecuaciones van enumeradas para facilitar la correción en caso de que haya algún typo (en la versión final solo las importantes deberían tener numeración).
    \item Por último, en el documento hay al menos tres problemas que arreglar: la notación no es totalmente consistente en todas las partes ya que he ido escribiendo en varios meses (arreglaré esto en la última lectura); del mismo modo, en algunas partes hay cosas técnicas que simplificar (algunos teoremas están enunciados sobre espacios polacos, la idea es no ir más allá de $\R^d$). Por otro lado, en algunas partes falta hacer una transición más suave entre los párrafos y/o secciones ya que he movido fragmentos de texto de un lado a otro.
\end{itemize}

\textbf{disclaimer:} la mayor parte de las imágenes las saqué de los papers que fui leyendo, mientras que las imágenes que no están citadas las generé usando ChatGPT (le pedí que me entregara los códigos para los gráficos). Por otra parte, nada de lo que está escrito en este documento fue generado por ChatGPT y tampoco los códigos adjuntos (por completitud, en el trabajo van implementaciones limpias y minimales de un VAE, de un modelo de difusión (con y sin SDE) y de una U-Net. También hay algunos códigos cortos del algoritmo de Sinkhorn, del puente browniano y otras cosas menores).